\section{Discussion and Limitations}
\label{sec:limits}
\subsection{What the Integration Model Establishes}
The useful unit of reuse is the association between a deployed member, a node capability, and a production run. It lets the same task use engineering quantities across different drivers while leaving execution and compensation attributable. The two-scenario tests support reconfiguration within the existing driver and scenario vocabulary; they do not measure the engineering effort to author a new scenario, validate its recipe windows, or integrate a new protocol. Likewise, fourteen registered engines describe an interface inventory, not fourteen independently evaluated controllers.

Separating proposal, admission, and disposition prevents several misleading success criteria. A recorded \texttt{keep} is the agent's judgment, a matching readback is a driver-level observation, and sustained process convergence is a separate experimental outcome. The process-freeze drill demonstrates why these layers cannot be collapsed. The LLM case studies further show that observability and scoring must be checked before decision quality can be compared: neither a task-completion string nor a single simulator-truth sample repairs a missing process window.

\subsection{Implementation and Evidence Boundaries}
The current design does not make asynchronous device I/O atomic with concurrent recipe changes. Recipe application and hold heartbeats also have different checks from agent-tool calls, and multi-node recovery can partly fail. Journals carry attribution but provide no claimed cryptographic integrity or guaranteed durable-before-actuation commit. The small-step policy is guidance, not a hard limiter. Cross-agent in-window oscillations, shared-node contention over longer horizons, and availability during storage failure require further evaluation. Batch-gated acquisition also leaves a monitoring gap between runs. These limits constrain the claim to the described supervisory paths.

Approval reduces permitted automation but does not establish comprehensive endpoint security, protect against a privileged compromised harness, or validate the approved engineering intent. Least-privilege deployment and industrial network separation remain separate responsibilities \cite{iec62443,nist80082}; the implementation is not a certified safety instrumented function \cite{iec61511}. The present evidence includes no physical hardware campaign, operator workload study, cross-engine comparison, or statistical equivalence test. Missing historical raw files and the integrated comparison-tool defect limit reproduction claims as detailed in Sec.~\ref{sec:repro}. Repairing these evidence paths, aligning process observations with simulator truth, and evaluating concurrent configuration changes are prerequisites to stronger deployment claims, not results of this paper.

\section{Conclusion}
\label{sec:conclusion}
\system{} integrates multi-agent coordination with heterogeneous industrial supervision through explicit node capabilities and shared production context. The architecture explains how connection, team deployment, acquisition, source-dependent execution checks, recipe operations, and human decisions cooperate without treating them as one atomic control function. Existing simulation and protocol experiments demonstrate the tested integration paths; fixed-controller trajectories and exploratory LLM runs answer different questions and are reported separately. The framework provides an inspectable foundation for supervised industrial agent applications, while physical validation, stronger execution consistency, and controlled agent comparisons remain open. Platform and benchmark sources are distributed under the repository's stated license; the simulator is maintained separately.\footnote{Artifact repositories: \url{https://github.com/kingdol666/AgentWorkShop} and \url{https://github.com/kingdol666/plc-node-simulator}. A frozen anonymous artifact and archival identifier must be prepared for the applicable submission stage.}

\section*{Disclosure}
AI-based assistants supported manuscript editing, code-assisted evidence checking, and figure preparation under author direction. The authors remain responsible for reviewing the claims, data, and final submission.
